\documentclass[11pt]{article}

\usepackage{amsmath,amsthm,amssymb}
\usepackage[margin=1.2in]{geometry}
\usepackage{hyperref}

\newtheorem{theorem}{Theorem}
\newtheorem{corollary}[theorem]{Corollary}
\newtheorem{lemma}[theorem]{Lemma}
\newtheorem{proposition}[theorem]{Proposition}
\theoremstyle{definition}
\newtheorem{definition}[theorem]{Definition}
\theoremstyle{remark}
\newtheorem{remark}[theorem]{Remark}

\title{Eigenvalue Positivity for the Hecke Staircase Element\\via Kazhdan--Lusztig Positivity}
\author{Clio}
\date{April 17, 2026 (v2)}

\begin{document}
\maketitle

\begin{abstract}
We prove that all nonzero eigenvalues of the Hecke staircase element
$\Pi_q = \prod_{j=1}^m (1 + T_{i_j})$ acting on any Specht module $V_\lambda$
are positive for all $q > 0$. The proof combines the Elias--Williamson positivity
theorem for Kazhdan--Lusztig basis structure constants with an exterior power
trace argument. As a corollary, the rank of $\Pi_q$ on $V_\lambda$ is constant
for $q > 0$. This note is intended to be self-contained: we include the
necessary definitions of the Iwahori--Hecke algebra, Specht modules,
W-graphs, and the staircase element.
\end{abstract}

%% ================================================================
\section{The Iwahori--Hecke algebra}\label{sec:hecke}
%% ================================================================

Let $n \geq 2$ and let $q$ be a formal parameter (or a positive real number).

\begin{definition}[Iwahori--Hecke algebra]
The \emph{Iwahori--Hecke algebra} $H_q(S_n)$ is the associative
$\mathbb{Z}[q,q^{-1}]$-algebra with generators $T_1, \ldots, T_{n-1}$
subject to the relations:
\begin{enumerate}
\item[(H1)] \emph{Quadratic relation:}\quad
  $T_i^2 = (q - 1)\,T_i + q$ \qquad for $1 \leq i \leq n-1$.
\item[(H2)] \emph{Braid relations:}\quad
  $T_i\,T_{i+1}\,T_i = T_{i+1}\,T_i\,T_{i+1}$ \qquad for $1 \leq i \leq n-2$.
\item[(H3)] \emph{Far commutativity:}\quad
  $T_i\,T_j = T_j\,T_i$ \qquad for $|i - j| \geq 2$.
\end{enumerate}
At $q = 1$, the quadratic relation becomes $T_i^2 = 1$, so
$H_1(S_n) \cong \mathbb{Z}[S_n]$ via $T_i \mapsto s_i$
(the simple transposition $(i,\,i{+}1)$).
For general~$q$, $H_q(S_n)$ is a free $\mathbb{Z}[q,q^{-1}]$-module
of rank~$n!$, with basis $\{T_w : w \in S_n\}$ indexed by
$T_w = T_{i_1} \cdots T_{i_\ell}$ for any reduced expression
$w = s_{i_1} \cdots s_{i_\ell}$.
\end{definition}

\begin{remark}
We will also use the parameter $v$ with $v^2 = q$.
The \emph{Kazhdan--Lusztig generator} is $C'_s = v^{-1}(1 + T_s)$
for a simple reflection~$s$.
This normalization is chosen so that the KL basis elements
$\{C'_w : w \in S_n\}$ have positivity properties.
\end{remark}

%% ================================================================
\section{Specht modules}\label{sec:specht}
%% ================================================================

\begin{definition}[Specht modules]
A \emph{partition} $\lambda = (\lambda_1 \geq \lambda_2 \geq \cdots
\geq \lambda_\ell > 0)$ of $n$ (written $\lambda \vdash n$)
indexes a \emph{Young diagram}: an array of $n$ boxes with
$\lambda_i$ boxes in row~$i$.
A \emph{standard Young tableau} (SYT) of shape~$\lambda$ is a
filling of the Young diagram with $1, \ldots, n$, each appearing once,
such that entries increase along each row and down each column.

The \emph{Specht module} $V_\lambda$ is an irreducible representation
of $H_q(S_n)$ (over $\mathbb{Q}(q)$, or over $\mathbb{Q}$ when $q = 1$
gives the classical Specht module for~$S_n$).
Its dimension is $f^\lambda = |\{\text{SYT of shape } \lambda\}|$,
computed by the hook-length formula.
The Specht modules $\{V_\lambda : \lambda \vdash n\}$ form a complete
set of pairwise non-isomorphic irreducible $H_q(S_n)$-modules
(for generic~$q$, i.e.\ $q$ not a root of unity).
\end{definition}

At $q = 1$, these are the classical irreducible representations of~$S_n$.
For general~$q$, Specht modules arise as \emph{cell modules} of the
Kazhdan--Lusztig cell structure: each two-sided cell corresponds to a
partition, and the cell module is the Specht module~$V_\lambda$.

%% ================================================================
\section{W-graphs and positivity}\label{sec:wgraph}
%% ================================================================

The key input to our proof is the non-negativity of certain matrix
entries in the Kazhdan--Lusztig basis.
This is most naturally expressed through the theory of W-graphs.

\begin{definition}[W-graph {\cite{KL79}}]
A \emph{W-graph} for $H_q(S_n)$ is a triple $(V, I, \mu)$ where:
\begin{itemize}
\item $V$ is a finite set of vertices (forming a basis of the representation),
\item $I \colon V \to 2^{\{1,\ldots,n-1\}}$ assigns a \emph{descent set}
  $I(x) \subseteq \{1, \ldots, n-1\}$ to each vertex,
\item $\mu \colon V \times V \to \mathbb{Z}[v^{-1}]$ assigns edge weights,
\end{itemize}
such that the action of $T_i$ on the basis $\{e_x : x \in V\}$ is:
\[
  T_i \cdot e_x \;=\;
  \begin{cases}
    q \cdot e_x & \text{if } i \notin I(x), \\[4pt]
    -e_x + \displaystyle\sum_{\substack{y \in V \\ i \notin I(y)}}
      \mu(y, x)\, v \cdot e_y
    & \text{if } i \in I(x).
  \end{cases}
\]
\end{definition}

\begin{theorem}[Kazhdan--Lusztig {\cite{KL79}}, Elias--Williamson {\cite{EW14}}]
\label{thm:KLpos}
Every Specht module $V_\lambda$ admits a W-graph in which
all edge weights satisfy $\mu(y, x) \in \mathbb{Z}_{\geq 0}[v^{-1}]$.
More generally, the structure constants for the product of
Kazhdan--Lusztig basis elements satisfy
\[
  C'_u \cdot C'_v \;=\; \sum_{w} h^w_{u,v}(v)\, C'_w,
  \qquad h^w_{u,v} \in \mathbb{Z}_{\geq 0}[v, v^{-1}].
\]
The non-negativity of $\mu$-coefficients was conjectured by
Kazhdan--Lusztig~\cite{KL79} and proved by
Elias--Williamson~\cite{EW14}.
\end{theorem}

The crucial consequence for us is:

\begin{proposition}\label{prop:1pTpos}
In the W-graph basis of $V_\lambda$, the matrix of $1 + T_i$ has
all entries in $\mathbb{Z}_{\geq 0}[v]$. Specifically:
\[
  (1 + T_i) \cdot e_x \;=\;
  \begin{cases}
    (1 + q)\, e_x & \text{if } i \notin I(x), \\[4pt]
    \displaystyle\sum_{\substack{y \in V \\ i \notin I(y)}}
      \mu(y, x)\, v \cdot e_y
    & \text{if } i \in I(x).
  \end{cases}
\]
Since $\mu(y,x) \in \mathbb{Z}_{\geq 0}[v^{-1}]$, the product
$\mu(y,x) \cdot v \in \mathbb{Z}_{\geq 0}[v^{-1}] \cdot v
\subseteq \mathbb{Z}_{\geq 0}[v, v^{-1}]$.
In fact, $1 + T_i = v \cdot C'_{s_i}$ and the entries lie in
$\mathbb{Z}_{\geq 0}[v]$ (no negative powers of~$v$ survive
because $C'_{s_i}$ acts on the cell module with Laurent polynomial
entries in $\mathbb{Z}_{\geq 0}[v^{-1}]$ and the extra factor of~$v$
shifts all exponents up by~$1$).
\end{proposition}

%% ================================================================
\section{The staircase element}\label{sec:staircase}
%% ================================================================

The long element $w_0 \in S_n$ is the unique permutation of maximal
length $m = \binom{n}{2}$.  A \emph{reduced word} for $w_0$ is a
sequence $(i_1, \ldots, i_m)$ with
$w_0 = s_{i_1} \cdots s_{i_m}$.

\begin{definition}[Staircase reduced word]
The \emph{staircase reduced word} for $w_0$ is:
\[
  \underbrace{s_1}_{k=1},\;\;
  \underbrace{s_2,\, s_1}_{k=2},\;\;
  \underbrace{s_3,\, s_2,\, s_1}_{k=3},\;\;
  \ldots,\;\;
  \underbrace{s_{n-1},\, s_{n-2},\, \ldots,\, s_1}_{k=n-1}.
\]
The $k$-th block is $(s_k,\, s_{k-1},\, \ldots,\, s_1)$, which
inserts element $k+1$ into positions $1$ through $k+1$.
This is a reduced word for~$w_0$ because it sorts
$(1, 2, \ldots, n)$ to $(n, n{-}1, \ldots, 1)$ by successively
moving each new largest element to the front.
\end{definition}

\begin{definition}[Hecke staircase element]
The \emph{Hecke staircase element} is:
\[
  \Pi_q \;=\; \prod_{k=1}^{n-1} \prod_{j=k}^{1} (1 + T_j)
  \;=\; \prod_{j=1}^{m} (1 + T_{i_j})
  \;=\; v^m \prod_{j=1}^{m} C'_{s_{i_j}},
\]
where $(i_1, \ldots, i_m)$ is the staircase reduced word.
\end{definition}

\begin{remark}[Why this reduced word?]
The staircase word has several special properties:
\begin{enumerate}
\item Its reversal is also a reduced word for~$w_0$
  (the ``anti-staircase''), which implies $f_w = f_{w^{-1}}$
  in the $q = 1$ group algebra expansion.
\item The block structure
  $(1 + T_k)(1 + T_{k-1}) \cdots (1 + T_1)$ at step~$k$
  is the divided symmetrization operator for the parabolic
  subgroup $S_{k+1}$.
  This gives $\Pi_q$ a cascade structure: it is a composition
  of partial symmetrizers of increasing rank.
\item At $q = 0$, the factors $(1 + T_i)$ become sorting operators
  $(1 + s_i)$, and $\Pi_0$ implements the well-known
  ``staircase sorting'' of permutations.
\item At $q = 1$, $\Pi_1 = \prod (1 + s_{i_j}) \in \mathbb{Q}[S_n]$,
  and its image on each Specht module $V_\lambda$ is the
  $H$-invariant subspace $V_\lambda^H$ for the parabolic
  subgroup $H = \langle s_1, s_3, s_5, \ldots \rangle$.
\end{enumerate}
\end{remark}

%% ================================================================
\section{Main results}
%% ================================================================

\begin{theorem}[Eigenvalue Positivity]\label{thm:main}
For every $\lambda \vdash n$ and $q > 0$, all nonzero eigenvalues of\/
$\Pi_q$ on the Specht module $V_\lambda$ are strictly positive.
\end{theorem}

\begin{corollary}[Rank Constancy]\label{cor:rank}
$\operatorname{rank}(\Pi_q|_{V_\lambda})$ is constant for all $q > 0$.
\end{corollary}

%% ================================================================
\section{Proof of trace non-negativity}
%% ================================================================

\begin{theorem}[Trace Non-negativity]\label{thm:trace}
For every partition $\mu \vdash n$, the trace $\operatorname{tr}_\mu(\Pi_q)$
of\/ $\Pi_q$ on the Specht module $V_\mu$ is a polynomial in $q$ with
non-negative integer coefficients.
\end{theorem}

\begin{proof}
\emph{Step~1 (Product expansion via Elias--Williamson).}
By Theorem~\ref{thm:KLpos}, the product of
Kazhdan--Lusztig basis elements has non-negative expansion:
\[
  \prod_{j=1}^m C'_{s_{i_j}}
  \;=\; \sum_{w \in S_n} a_w(v)\, C'_w
  \qquad\text{with } a_w \in \mathbb{Z}_{\geq 0}[v].
\]
This follows because the product $C'_{s_{i_1}} \cdot C'_{s_{i_2}}
\cdots C'_{s_{i_m}}$ is built up one factor at a time, and at each
step, multiplying by $C'_{s_{i_j}}$ preserves non-negativity of the
expansion coefficients by the Elias--Williamson positivity of the
structure constants~$h^w_{u,v}$.

\emph{Step~2 (Trace of a KL basis element on a cell module).}
Let $V_\mu$ be the cell module (Specht module) corresponding to~$\mu$.
In the KL cellular basis $\{\bar{C}_x : x \in \mathrm{cell}(\mu)\}$,
left multiplication by $C'_w$ gives:
\[
  C'_w \cdot \bar{C}_x
  \;=\; \sum_{y \in \mathrm{cell}(\mu)} h_{w,x,y}(v)\, \bar{C}_y
\]
where $h_{w,x,y} \in \mathbb{Z}_{\geq 0}[v, v^{-1}]$ by KL positivity.
(For cell modules, in fact $h_{w,x,y} \in \mathbb{Z}_{\geq 0}[v^{-1}]$.)
The trace is therefore:
\[
  \operatorname{tr}_\mu(C'_w)
  \;=\; \sum_{x \in \mathrm{cell}(\mu)} h_{w,x,x}(v)
  \;\in\; \mathbb{Z}_{\geq 0}[v, v^{-1}].
\]

\emph{Step~3 (Combining and specializing).}
Since $\Pi_q = v^m \prod_{j} C'_{s_{i_j}}$:
\[
  \operatorname{tr}_\mu(\Pi_q)
  \;=\; v^m \sum_{w} a_w(v) \cdot \operatorname{tr}_\mu(C'_w)
  \;\in\; \mathbb{Z}_{\geq 0}[v, v^{-1}].
\]
Now, $\operatorname{tr}_\mu(\Pi_q)$ lies in $\mathbb{Z}[q] = \mathbb{Z}[v^2]$
because it is the trace of a $\mathbb{Z}[q]$-matrix on a $\mathbb{Z}[q]$-module
(the representation matrices in the $T$-basis have entries in $\mathbb{Z}[q]$).
So all odd powers of~$v$ vanish.
A polynomial in $\mathbb{Z}_{\geq 0}[v, v^{-1}] \cap \mathbb{Z}[v^2]$
has non-negative coefficients when re-expressed in $q = v^2$.
Hence $\operatorname{tr}_\mu(\Pi_q) \in \mathbb{Z}_{\geq 0}[q]$.
\end{proof}

%% ================================================================
\section{Proof of eigenvalue positivity}
%% ================================================================

\begin{proof}[Proof of Theorem~\ref{thm:main}]
\emph{Step~1 (Elementary symmetric functions of eigenvalues).}
Let $r_\lambda$ denote the rank of $\Pi_q$ on $V_\lambda$ over
$\mathbb{Q}(q)$, and let $\alpha_1(q), \ldots, \alpha_{r_\lambda}(q)$
be the nonzero eigenvalues (as algebraic functions of~$q$).
For $1 \leq k \leq r_\lambda$, let
\[
  e_k(q) \;=\; \sum_{1 \leq j_1 < \cdots < j_k \leq r_\lambda}
    \alpha_{j_1}(q) \cdots \alpha_{j_k}(q)
\]
be the $k$-th elementary symmetric function of the nonzero eigenvalues.
By the standard identity relating exterior powers to elementary
symmetric functions:
\[
  e_k(q) \;=\; \operatorname{tr}\!\Bigl(
    \textstyle\bigwedge^k(\Pi_q|_{V_\lambda})
    \text{ on } \bigwedge^k V_\lambda
  \Bigr).
\]
This uses the multiplicativity of exterior powers:
$\bigwedge^k(AB) = (\bigwedge^k A)(\bigwedge^k B)$, so
the product $\Pi_q = \prod(1 + T_{i_j})$ lifts to
$\bigwedge^k(\Pi_q) = \prod \bigwedge^k(1 + T_{i_j})$.

\emph{Step~2 (Decomposition into Specht modules).}
As an $H_q(S_n)$-module,
\[
  \textstyle\bigwedge^k V_\lambda
  \;\cong\; \bigoplus_\mu V_\mu^{\oplus a_{k,\mu}}
\]
with $a_{k,\mu} \geq 0$.
By Theorem~\ref{thm:trace}:
\[
  e_k(q) \;=\; \sum_\mu a_{k,\mu} \cdot \operatorname{tr}_\mu(\Pi_q)
  \;\in\; \mathbb{Z}_{\geq 0}[q].
\]

\emph{Step~3 (Non-vanishing).}
Over $\mathbb{Q}(q)$, the matrix $\Pi_q|_{V_\lambda}$ has rank~$r_\lambda$.
Hence $e_{r_\lambda} = \alpha_1 \cdots \alpha_{r_\lambda} \neq 0$ as a
rational function, and therefore as a polynomial in~$q$
(since $e_k \in \mathbb{Z}[q]$).  More generally, $e_k \neq 0$
for $1 \leq k \leq r_\lambda$.
A nonzero polynomial in $\mathbb{Z}_{\geq 0}[q]$ is strictly positive
for all $q > 0$.

\emph{Step~4 (Positivity of all eigenvalues).}
The nonzero eigenvalues $\alpha_1, \ldots, \alpha_{r_\lambda}$ of
$\Pi_q$ on $V_\lambda$ are real for $q > 0$ (as eigenvalues of a
real matrix when $q$ is specialized to a positive real number).
They are the roots of the reduced characteristic polynomial:
\[
  p(x) \;=\; \sum_{k=0}^{r_\lambda} (-1)^k\, e_k(q)\, x^{r_\lambda - k},
  \qquad e_0 = 1.
\]
We claim $p(x) \neq 0$ for all $x \leq 0$ (when $q > 0$).
For $x \leq 0$, consider consecutive terms:
\begin{itemize}
\item $(-1)^k\, e_k(q)\, x^{r_\lambda - k}$
  has sign $(-1)^k \cdot (+) \cdot (-1)^{r_\lambda - k}
  = (-1)^{r_\lambda}$.
\item $(-1)^{k+1}\, e_{k+1}(q)\, x^{r_\lambda - k - 1}$
  has sign $(-1)^{k+1} \cdot (+) \cdot (-1)^{r_\lambda - k - 1}
  = (-1)^{r_\lambda}$.
\end{itemize}
All terms have the common sign $(-1)^{r_\lambda}$, so
$p(x) \neq 0$ for $x \leq 0$.
Therefore all roots of $p$ are positive.
\end{proof}

\begin{proof}[Proof of Corollary~\ref{cor:rank}]
Since all $r_\lambda$ eigenvalues are strictly positive for $q > 0$,
none can vanish, so the rank cannot drop below~$r_\lambda$.
And rank cannot exceed~$r_\lambda$ (the generic rank over $\mathbb{Q}(q)$).
\end{proof}

%% ================================================================
\section{Computational verification}\label{sec:comp}
%% ================================================================

The trace non-negativity and eigenvalue positivity have been verified
independently by explicit computation of eigenvalue polynomials for
all $\lambda \vdash n$ with $n \leq 8$.

\paragraph{Method.}
For each $\lambda \vdash n$, we construct the Specht module $V_\lambda$
using seminormal matrix units, compute the matrix of each factor
$(1 + T_{i_j})$, multiply them to obtain $\Pi_q|_{V_\lambda}$,
and compute its characteristic polynomial in $\mathbb{Z}[q][x]$.
The nonzero eigenvalue polynomials are then extracted.

\paragraph{Sample results.}
\begin{center}
\begin{tabular}{cccc}
\hline
$\lambda$ & $\dim V_\lambda$ & rank of $\Pi_q$ & eigenvalue polynomials \\
\hline
$(3)$ & 1 & 1 & $(1+q)^3$ \\
$(2,1)$ & 2 & 1 & $q(1+q)$ \\
$(1,1,1)$ & 1 & 1 & $q^3$ \\
$(4)$ & 1 & 1 & $(1+q)^6$ \\
$(3,1)$ & 3 & 2 & $q^2(1+q)^2(1+q+q^2),\; q(1+q)^4$ \\
$(2,2)$ & 2 & 1 & $q^2(1+q)^2(2q+1)$ \\
$(2,1,1)$ & 3 & 2 & $q^3(1+q)^2(1+q+q^2),\; q^4(1+q)^2$ \\
$(1^4)$ & 1 & 1 & $q^6$ \\
\hline
\end{tabular}
\end{center}

\noindent
All listed eigenvalue polynomials have non-negative coefficients
and are strictly positive for $q > 0$, confirming
Theorems~\ref{thm:trace} and~\ref{thm:main} in these cases.
Verification extends through $n = 8$ ($|\{\lambda \vdash 8\}| = 22$
partitions, all confirmed).

\paragraph{Independence from reduced word.}
Trace non-negativity holds for the staircase word specifically.
Other reduced words for $w_0$ can yield elements $\Pi'_q$ whose
traces have negative coefficients.  The staircase word is
distinguished by its cascade (parabolic) structure.

%% ================================================================
\section{Remarks}
%% ================================================================

\begin{remark}[Minimal hypotheses]
The proof requires \emph{only}:
(i) the identification $1 + T_s = v \cdot C'_s$,
(ii) Elias--Williamson positivity of KL structure constants, and
(iii) the decomposition of exterior powers into Specht modules.
No case analysis, no explicit eigenvalue computation, no restriction
on~$n$, and no assumption on the partition~$\lambda$.
\end{remark}

\begin{remark}[Connection to the $q = 0$ limit]
At $q = 0$, the Hecke algebra degenerates to the $0$-Hecke algebra
$H_0(S_n)$, where $T_i^2 = -T_i$.  The factors $(1 + T_i)$ become
idempotent: $(1 + T_i)^2 = (1 + T_i)$.  The staircase element
$\Pi_0$ acts as the ``staircase sorting operator,'' a projection
whose image on~$V_\lambda$ is the $H$-invariant subspace
$V_\lambda^H$.  Theorem~\ref{thm:main} shows that the positivity
persists for all $q > 0$, interpolating smoothly from this $q = 0$
projection.
\end{remark}

\begin{remark}[Application]
The eigenvalue positivity theorem is the key input to the
``Frobenius squeeze'' proof of the $H$-invariant theorem
$\operatorname{rank}(\Pi_q|_{V_\lambda}) = \dim(V_\lambda^H)$:
it provides the per-irreducible lower bound that, combined with
the total rank upper bound from Frobenius reciprocity, forces
equality in every component.
\end{remark}

\begin{thebibliography}{EW14}

\bibitem[EW14]{EW14}
B.~Elias and G.~Williamson,
\emph{The Hodge theory of Soergel bimodules},
Ann.\ of Math.\ \textbf{180} (2014), 1089--1136.

\bibitem[KL79]{KL79}
D.~Kazhdan and G.~Lusztig,
\emph{Representations of Coxeter groups and Hecke algebras},
Invent.\ Math.\ \textbf{53} (1979), 165--184.

\bibitem[Hum90]{Hum90}
J.~E.~Humphreys,
\emph{Reflection Groups and Coxeter Groups},
Cambridge Univ.\ Press, 1990.

\bibitem[Mat99]{Mat99}
A.~Mathas,
\emph{Iwahori--Hecke Algebras and Schur Algebras of the Symmetric Group},
Univ.\ Lecture Series~\textbf{15}, AMS, 1999.

\end{thebibliography}

\end{document}
